% ============================================================
%  Harald Høffding, Den menneskelige Tanke
%  Dens Former og dens Opgaver
%  København og Kristiania: Gyldendalske Boghandel, Nordisk Forlag, 1910.
%  TRANSCRIPTION (Danish)
%  Source: KB scan (scan.pdf). Page-offset: PDF = printed + 15.
% ============================================================
\documentclass[12pt,a4paper]{book}
\usepackage[utf8]{inputenc}
\usepackage[T1]{fontenc}
\usepackage{libertinus}
\usepackage{libertinust1math}
\usepackage{textalpha}   % Greek words typed directly
\usepackage{graphicx}    % for \reflectbox (the old Danish »ɔ:« id-est mark)
\DeclareUnicodeCharacter{0254}{\reflectbox{c}}  % ɔ (U+0254) = reversed c
\usepackage[danish]{babel}
\usepackage[protrusion=true,expansion=true]{microtype}
\usepackage{setspace}
\usepackage[margin=1.2in]{geometry}
\usepackage{enumitem}
\usepackage{fancyhdr}
\setlength{\headheight}{28pt}
\addtolength{\topmargin}{-13.5pt}
\usepackage{hyperref}

\hypersetup{
  pdftitle={Den menneskelige Tanke},
  pdfauthor={Harald Høffding},
  hidelinks
}

\pagestyle{fancy}
\fancyhf{}
\fancyhead[LE,RO]{\thepage}
\fancyhead[RE]{\textit{Den menneskelige Tanke}}
\fancyhead[LO]{\textit{\leftmark}}

\setstretch{1.3}

\title{\textbf{Den menneskelige Tanke}\\[6pt]
  \large Dens Former og dens Opgaver}
\author{Harald Høffding}
\date{København og Kristiania:\\ Gyldendalske Boghandel, Nordisk Forlag, 1910}

\begin{document}
\maketitle
\thispagestyle{empty}
\newpage

\tableofcontents
\newpage

%% ============================================================
%%  FORORD (foreword, August 1910)
%% ============================================================
\chapter*{Forord}
\addcontentsline{toc}{chapter}{Forord}
\label{ch:forord}

\noindent Tankelivet, som det er min Bogs Opgave at skildre, er
inderligt sammenvævet med Aandslivets andre Sider. Jeg har søgt at
vise, at Tanken ikke kan tjene Aandslivet idethele, naar den ikke
først og fremmest bliver Herre paa sit eget Omraade og udfolder sig
efter sine egne Love som en selvstændig Side af Aandslivet. At denne
Udfoldelse, naar den ret gennemføres, vil medføre store og indgribende
Ændringer i det personlige Liv idethele, derom bliver jeg mere og mere
overbevist. Der er her Modsætninger og Problemer, som ville antage
stedse skarpere Former, og jeg føler ingen Beundring for dem, der her
ikke føle nogen Brod, skønt de have tilstrækkelig Adgang til at lære
Tankelivet at kende.

Som jeg idethele har betonet den personlige Faktor i al Tænken, jo
mere denne nærmer sig Grænseomraaderne, saaledes har jeg heller ikke
lagt Skjul paa, at min egen Holdning paa Tankens sidste Station har en
personlig Karakter, og jeg har ikke kunnet undlade at udtale,
hvorledes man efter min Erfaring er tilmode paa denne Station.

Jeg har endnu Meget at lære. Udfyldning og Berigtigelse vil mit
Arbejde trænge til paa mange Punkter. Men ud over den Ramme, jeg her
har givet, kommer jeg ikke. Og den afgiver for mig en god
Arbejdsplads. Maatte der være En og Anden, som ogsaa kunde finde
Arbejde her, selvom dette Arbejde skulde føre ham til at ændre baade
ved Rammen og ved dens Indhold.

\bigskip
\noindent August 1910.\hfill \emph{Harald Høffding.}

%% ============================================================
%%  I. TANKENS PSYKOLOGI (FUNKTIONER)   [pp. 1--108]
%% ============================================================
\chapter*{I. Tankens Psykologi (Funktioner)}
\addcontentsline{toc}{chapter}{I. Tankens Psykologi (Funktioner)}
\markboth{I. Tankens Psykologi}{I. Tankens Psykologi}
\label{ch:psykologi}

\begin{center}---\end{center}

\section*{A. Psykisk Energi}
\addcontentsline{toc}{section}{A. Psykisk Energi}
\label{sec:psykisk-energi}

%% ---------- printed pp. 1--10 ----------
\noindent 1. Den Opgave, jeg har sat mig i denne Bog, er at undersøge
den menneskelige Tanke --- en Opgave, alle Filosofer have sat sig,
hvad enten de have været sig det bevidst eller ikke. Ti alle
Spørgsmaal, der kunne opkastes angaaende Livet og Tilværelsen, afhænge
--- baade med Hensyn til den Maade, paa hvilken de stilles, og med
Hensyn til den Maade, paa hvilken de besvares --- af den Opfattelse,
man har af Tankens Natur og Virkemaade. Det er da ogsaa et Omrids af
min Filosofi, jeg her vil give. Det vil sige: jeg vil give en
Fremstilling af de Tanker, til hvilke jeg atter og atter er kommen
tilbage, naar jeg har syslet med saadanne Spørgsmaal som: Videnskabens
Forudsætninger, Betydningen af Videnskabens Resultater for
Livsanskuelsen og Forholdet mellem videnskabelig Erkendelse og andre
Sider af det menneskelige Aandsliv.

2. Vi kende Tanken bedst som Eftertanke. Ordet Eftertanke bruger jeg
ikke blot i den egentlige Betydning „i sin Tænken at \emph{søge efter}
Noget eller \emph{forfølge} Noget“, men ogsaa i den Betydning, at
Tanken \emph{kommer bagefter} en umiddelbar Oplevelse. Eftertanken
forudsætter da, at vi have set, følt, oplevet Noget. Under Oplevelsen
gik vi maaske helt op i det; men nu søge vi at fremdrage det paany
enten for at dvæle ved det Oplevede igen, eller for at gennemgaa dets
forskellige Dele eller Egenskaber eller dets Forhold til andre
Oplevelser. I den blotte Dvælen er Eftertanken Et med en Skuen, en
blot Forestillen. Men en saadan Dvælen vil neppe kunne vedvare længe i
fuldstændigt hvilende Form. Uvilkaarligt vil der --- maaske netop paa
Grund af den Interesse, det Oplevede har, udløses en Tankebevægelse,
der tjener til at fremhæve dets indre og ydre Forhold. Saaledes naar
Digteren uvilkaarligt griber til Lignelser og Billeder for ret at lade
sit Emne komme til Udtryk. Eller naar maaske netop Kontrasten til
andre Oplevelser virker til Belysning af den enkelte Oplevelses
ejendommelige Karakter. Baade Lighed og Forskel kunne saaledes virke
til at kaste et klarere Lys over, hvad der var givet i den oprindelige
Oplevelse, og Eftertanken staar derfor ikke nødvendigvis i Modsætning
til denne, men kan virke i dens Tjeneste og i dens Aand, naar den selv
er forsvunden.

Vi tænke ved Ordet Oplevelse nærmest paa Tilstande, ved hvilke vi
væsentligst vare modtagende. Men Eftertanken kan ogsaa tjene det
uvilkaarlige Aandsliv, hvor dette fremtræder som Trang eller Stræben.
Paa instinktiv eller genial Maade føres den uvilkaarlige Stræben ofte
til Maalet, uden at nogen Eftertanke behøves, og maaske bedre end det
kunde ske ved nogen Eftertankes Hjælp. Naar Eftertanken kommer til at
virke i vor Handlen, vil det være, fordi den uvilkaarlige Stræben
fandt Modstand. Der kan da behøves en Besindelse over, hvad det
egentligt er, der stræbes efter, over dette Maals Forhold til andre
mulige eller virkelige Bestræbelser og over de Veje og Midler, der
kunne føre til Maalet. Ogsaa her er det gennem Fremdragen af Ligheder
og Forskelle, at Eftertanken yder sin Tjeneste. Og ogsaa her kan der
fremtræde en Forskel mellem Dvælen og Bevægelse. Det kan være, at
Forestillingen om Maalet holdes fast uden store Svingninger, og at
Forholdet til Veje og Midler eller til andre Muligheder spiller en
underordnet Rolle, ligesom naar Styrmanden stirrer paa Kompasset eller
Sømærket, og Haanden uvilkaarligt bevæger Roret efter hvad han ser.
Men det kan ogsaa være, at der foregaar og maa foregaa den stærkeste
Bevægelse frem og tilbage mellem forskellige Muligheder, inden den
afgjørende Koncentration, der kaldes Beslutning, indtræder.

Eftertankens Opstaaen er som en Slags Opvaagnen, ikke fra Søvn eller
fra Ubevidsthed, men fra det uvilkaarlige Sjæleliv og dets Oplevelser,
fra Tilstande, der foreløbigt kunne karakteriseres som Sansning,
Anskuen, Forestillen, Stemning, Sindsbevægelse, Trang og Stræben. I
sin dvælende Form staar Eftertanken nærmest ved det uvilkaarlige
Sjæleliv. I sin Stræben efter at fremdrage Ligheder og Forskelle
fjerner den sig mere fra det Uvilkaarlige. Sprogbrugen betegner ret
træffende denne Forskel ved at skelne mellem at \emph{tænke paa} Noget
og at \emph{tænke over} Noget. Det er den sidste Art af Eftertanke,
den diskursive, bevægelige, vi i det Følgende mest faa Brug for, da
det er den, vi bedst kunne iagttage, og tillige den, der spiller den
afgørende Rolle ved Tankens Opgaver. Kun for denne Art Eftertanke
existerer der Problemer. I det Følgende forstaar jeg derfor ved
Tænkning denne Art af Eftertanke. Og dog er der kun en Gradsforskel
mellem de to Arter. Kun et kort Øjeblik er dvælende Forestillen,
Fastholden af et Erindringsbillede mulig; nøjagtig Iagttagelse viser,
at vi i større eller mindre Rytmer gaa til og fra det Emne, vi tænke
paa. Og enhver Tilbagevenden kan medføre en Genkendelse, saa at
Lighedsforholdet kommer til at spille en Rolle ligesom ved den
diskursive Eftertanke. Maaske var Forestillingen os ny i første
Øjeblik; men gennem den gentagne Dvælen bliver den kendt. Ogsaa afset
herfra udelukker den dvælende Eftertanke ingenlunde, at mangehaande
andre Forestillinger og Erfaringer kunne medvirke og maaske netop
betinge den Interesse, der gør fortsat Fastholden af en enkelt
Forestilling mulig. Hvad der gælder for den diskursive Eftertanke, vil
derfor ogsaa til en vis Grad gælde for den dvælende.

Det fremgaar af denne Beskrivelse, at Tankens Arbejde bestaar i
Sammenligning, i at skelne eller at finde Ligheder. Naar jeg tænker
over Noget, søger jeg at bestemme dets Forhold til noget Andet, og jo
flere Ligheds- og Forskelsforhold, jeg finder, des bedre har jeg
gennemtænkt det. Det gælder naturligvis ogsaa, naar det er et og samme
Emne, jeg undersøger i forskellig Sammenhæng og under forskellige
Forhold; jeg sammenligner da de forskellige Maader, Emnet optræder paa
i de forskellige Tilfælde.

Men Sammenligning forudsætter igen Sammenholden eller Sammenfatten,
Syntese. Jeg kan kun sammenligne to Emner, naar jeg enten har dem
givne samtidigt for mig i min Sansning, i min Erindring eller i min
Fantasi, eller ogsaa bevarer en Forestilling om det ene, medens det
andet foreligger for sanselig Iagttagelse. Hvis stadigt det ene Emne
ganske forsvandt, naar det andet dukkede op, vilde ingen Sammenligning
kunne finde Sted. Tankearbejdet bestaar i en Sammenfatten, ved hvilken
Ligheder og Forskelle blive bestemmende for de Forhold, i hvilke
Emnerne stilles til hverandre i den Sammenhæng, der bliver Resultatet
af Arbejdet. Syntesen er Udslag af Energi, af Evne til at udføre et
Arbejde.

Modstanden, som overvindes gennem dette Arbejde, skyldes de Emner, mod
hvilke Eftertanken rettes, og som ere givne i det uvilkaarlige
Sjæleliv, --- i Sansefornemmelse, Erindring eller Fantasi, Lyst- og
Ulystfølelse, Trang og Stræben.

Det sammenfattende Arbejde maa blive des vanskeligere, jo større den
Mangfoldighed af Emner er, der skulle bringes i Forhold til hverandre.
Indholdsfylden kan være saa stor, at den psykiske Energi, der i det
enkelte Tilfælde er forhaanden, ikke kan magte den. Det truer da med
Opløsning af Sjælelivet, med Overgang til Kaos.

Dog kommer det ikke blot an paa Mangfoldigheden, men ogsaa paa den
Grad af Forskellighed, som Emnerne frembyde. Jo mere de ligne
hverandre, saa at de kunne sammensmelte, eller at dog det Arbejde, der
er gjort for det enes Vedkommende, ikke behøver at gentages for det
andets Vedkommende, des mindre Arbejde kræves der. Emner, der staa i
stærk indbyrdes Modsætning eller maaske endog Modstrid, kræve en stor
Energi, naar de alle skulle medoptages i Sjælelivets bevidste
Sammenhæng. Det gælder baade for Forskeren, der skal finde et
Synspunkt, ud fra hvilket det Stridende kan staa som Et, og for den
praktisk Stræbende, der skal underordne de stridende Tilskyndelser og
Interesser under Tanken om et stort og ledende Formaal.

Det vil for det Tredie komme an paa, hvor selvstændig en Plads de
enkelte Emner skulle indtage i den Sammenhæng, paa hvilken der
arbejdes. En skematisk Helhed er lettere at danne end en Helhed, der
lader de enkelte Emners Ejendommeligheder hver for sig komme til deres
Ret. Netop Ligheden kan her gøre en Vanskelighed, da den let fører til
Sammensmeltning og derved til Ophævelse af enkelte Emners
Selvstændighed. Paa den anden Side kan fra dette Synspunkt udpræget
Forskel eller Modsætning være en Lettelse.

Der kræves fremdeles et større Arbejde, jo fjernere Emnerne ligge
hverandre i Tid og Rum. Det Samtidige og det Nære slutter sig
uvilkaarligt sammen, og Sammenhængen paatvinger sig maaske endog. Og
forlade vi de foreliggende Emner for at fordybe os i Erindringens og
Fantasiens Verden, kan ogsaa her Sammenhængen ofte danne sig med
Lethed; „med stor Lethed fæste Tankerne Bo hos hverandre“, men naar
det gælder at sammenholde et nærværende Emne med fjerne Erindringer
eller med Idealer, der kun i enkelte begejstrede Øjeblikke have været
ret fremme i Sjælen, da stilles der særlige Fordringer til psykisk
Energi. I Øjeblikkets Situation at have sine Erfaringer, sin
Eftertankes og sin Idealitets Skatte til fuld Raadighed er maaske det
største Vidnesbyrd om sjælelig Energi.

Det spiller endeligt en Rolle, hvor nøje og inderligt en Forbindelse
Emnerne trods deres Forskellighed skulle indgaa. Der er her mange
Grader, ligefra helt udvortes Sammenstillen og Grupperen til den indre
Sammenhæng, som et videnskabeligt Begreb, et kunstnerisk udformet
Billede eller en stor Beslutning kan udtrykke.

Foruden disse fem Forhold --- der kunne siges at danne en Analogi til
hvad Massen er ved den fysiske Energi --- maa der ogsaa tages Hensyn
til den Hurtighed, hvormed Syntesen foregaar. Geniet og
Villiesmennesket kunne i et Nu anvende en Energikapital, som andre
Mennesker maa fordele paa et længere Tidsrum. Paa Følelseslivets
Omraade fremtræder denne Modsætning som en Modsætning mellem Heftighed
og Inderlighed, Sindsbevægelse og Sindelag\footnote{Se herom min
\emph{Psykologi} IV, 7; VI E, 4--5.}, og analoge Modsætninger gælde
for de andre Sider af Sjælelivet. I Forskningen danner Opdagelsen ofte
en Analogi til Sindsbevægelsen, Bevisførelsen til Sindelaget.

Den psykiske Energi, der ytrer sig i Syntesen, maale vi ved at tage
Hensyn til alle disse Omstændigheder. Om exakt Maaling er der her
naturligvis ikke Tale, kun om et vurderende Skøn. Men saasnart vi søge
at begrunde, hvorfor et Tankearbejde, en Aandsretning, en Livsanskuelse
kan kaldes lavere eller højere end en anden, maa vi tage Hensyn til
alle disse Omstændigheder, og vi gøre det i Virkeligheden ogsaa, mere
eller mindre bevidst, og mere eller mindre nøjagtigt. Der er mindre
psykisk Energi til Raadighed i Drømmetilstanden, hvor de skiftende
Emner udløse hver sit Forestillingsløb, end i den vaagne Tilstand,
hvor herskende Formaalstanker medføre en Hovedretning. Det er en rent
psykologisk Maalestok, vi her anvende. Og selvom vi ved Dannelsen af
Begrebet psykisk Energi have havt Vejledning i Analogien med den
fysiske Energis Begreb, saa er hint Begreb dog bygget paa sine
særegne Erfaringer, ved hvilke fysiske eller fysiologiske Iagttagelser
og Undersøgelser ikke behøve at spille nogen afgørende Rolle. Derfor
behøve vi ikke at miskende Betydningen af de Forsøg, der fra
experimental-psykologisk og fysiologisk Side ere gjorte paa at maale
det Hjernearbejde, der forudsættes ved visse elementære
Aandsvirksomheder, f.\ Ex. Regnen eller Husken, gennem dets Forhold
til det Muskelarbejde, der kan øves samtidigt dermed. Man kan mærke
saadant Tankearbejde ved den Forminskelse, det fremkalder i det
samtidigt udførte Muskelarbejde. Men selve denne Maaling forudsætter,
at man ad rent psykologisk Vej kan skelne mellem vanskeligere og
lettere Aandsarbejde; ellers bevæger man sig i en Kreds. Det
Hjernearbejde, der svarer til det i det Foregaaende beskrevne
aandelige Arbejde, kender man jo ikke direkte, kun gennem selve det
Aandsarbejde, til hvilket det antages at svare.

3. Eftertanken forudsætter Emner, paa hvilke den kan øve sin
sammenfattende og sammenlignende Virksomhed. Men nu disse Emner?
Gælder Syntesens Lov ogsaa for dem? Vi staa her ved et Spørgsmaal af
en ejendommelig Vanskelighed, der hænger nøje sammen med Sjælelivets
hele Natur. Det er først Eftertanken, der lærer os Emnerne at kende.
Den er det Medium, gennem hvilket de ere os tilgængelige. Det
uvilkaarlige Sjæleliv i dets umiddelbare Udfoldelse maa ikke uden
videre beskrives og bedømmes efter den Maade, paa hvilken det viser
sig for Eftertanken. Det er et Liv, indenfor hvilket vi ikke have Lov
til at forudsætte Egenskaber og Bestemmelser, der først blive til
gennem Eftertankens Arbejde. Naar vi i enkelte Øjeblikke kunne fordybe
os i dette uvilkaarlige, men derfor ikke ubevidste Liv saaledes, at vi
senere kunne bevare en Erindring om det, staar det for os som en
Strøm, en Hengliden af mangfoldige Emner, hvilke vi betegne som
Sansninger, Erindringer, Stemninger, Bestræbelser o.\ s.\ v. Det er
ikke et Kaos af selvstændige Emner; der er allerede i det uvilkaarlige
Sjæleliv Sammenhænge og Helheder saavel som Forskelligheder og
Modsætninger. Ethvert Ord, vi bruge om det, vi kunne fastholde fra de
uvilkaarlige Tilstande, vidner derom. Mest paafaldende er vel
Kontinuiteten; derfor frembyder Ordet „Strøm“ sig saa let om det, vi
kalde det umiddelbart Givne. Men det er ikke nogen ensformig Strøm.
Der gør sig Kvalitetsforskelle gældende; ellers kunde vi ikke skelne
mellem de Arter af Emner, der ovenfor kaldtes Sansninger, Erindringer
o.\ s.\ v. Og der kan være Hvirvler i Strømmen, eller Slyngninger, saa
at Retningen ikke holdes i lige Linie. Det vil sige, der kan gøre sig
Stræben gældende af forskellig Tendens, af mere eller mindre modsat
Retning. Dermed hænger det igen sammen, at der kan indtræde
Hæmninger, medens der i andre Tilfælde, hvor de forskellige Tendenser
gaa i samme Retning, kan fremtræde en Lettelse eller en Fremskynden.

Gennem de ubestemte og billedlige Udtryk, der ene staa til vor
Raadighed til at antyde det Grundlag af Emner, Eftertanken
forudsætter, skinner det igennem, at det ikke danner nogen absolut
Modsætning til, hvad der fremkommer som Følge af Eftertankens Arbejde.
Vi finde ogsaa hist baade Enhed og Mangfold, baade Kontinuitet og
Diskontinuitet, og vi finde mange Grader og Nuancer af disse Forhold.
Det er altsaa ikke berettiget med den ældre engelske Skole
(\emph{Locke}, \emph{Hume}, \emph{Mill}) at opfatte det umiddelbart
Givne i Sjælelivet paa atomistisk Maade, som om det bestod i en
Mangfoldighed af indbyrdes selvstændige Elementer. Baade paa det
psykiske og paa det fysiske Omraade er det først Eftertanken, der
fører til at danne Atombegreber, ledet af Trangen til at gøre
Forskellighederne saa tydelige som muligt, for at Forbindelsen af dem
kan blive saa klar og anskuelig som muligt. Den Enhed, hvorom
Eftertanken vidner, og som den selv søger at tilvejebringe, bliver
derved gaadefuld. Men det er ogsaa uberettiget, som i vore Dage især
\emph{William James}\footnote{\emph{The Stream of Thought}. Mind 1884
(genoptrykt som Kap.\ 9 i \emph{Principles of Psychology}). Senere
især \emph{A pluralistic universe} (1909), Sect.\ 7.} og \emph{Henri
Bergson}\footnote{\emph{Les données immédiates de la conscience}
(1888). Senere især i \emph{Introduction à la Métaphysique}. (Revue de
Métaphysique et de Morale 1903).} have forsøgt, at fremstille hint
uvilkaarlige Liv som staaende i den skarpeste Modsætning til
Eftertanken og dens Arbejde. Det bliver da til Syvende og Sidst det
Ubeskrivelige og Ubegribelige, og det vilde blive en uløselig Gaade,
hvorledes der overhovedet kan gives Kundskab om hint „umiddelbart
Givne“, af hvilket dog Eftertanken stadigt skal øse. James og Bergson
betegne en betydningsfuld Reaktion mod Intellektualismens
Misforstaaelse og Overvurdering af Eftertankens Betydning, en
Reaktion, der i flere Henseender minder om \emph{Hamann}’s,
\emph{Herder}’s og \emph{Jacobi}’s Optræden mod \emph{Kant}. Ikke den
rene Fornuft, men det aandelige Livs umiddelbare Strøm er, hævde de,
det Højeste og bringer Sandheden frem. Tankens Distinktioner og
Definitioner ere nødvendige, for at vi praktisk kunne finde os til
Rette i Livet. Og den analyserende Videnskab har ført Modsætningen til
det umiddelbart Givne videre, en Modsætning, allerede selve Sproget
begunstiger ved at danne selvstændige Ord for de enkelte Emner og ved
at lægge rumlige Forhold til Grund for vore Anskuelsesbilleder. Det
er, som om der er sket et Syndefald, da Eftertanken vaagnede og vendte
sig forskende mod det uvilkaarlige Liv, af hvilket den selv er
udsprungen.

Uden et Slægtskab mellem det uvilkaarlige Sjæleliv og Eftertanken
vilde det være umuligt at forstaa, hvorledes denne kunde opstaa. Men i
selve de Beskrivelser, James og Bergson give af, hvad der for dem i
skarp Modsætning til al Tankebearbejdelse er det umiddelbart Givne,
--- Beskrivelser, der ere Mesterstykker af deskriptiv Psykologi, ---
angives det dog udtrykkeligt, at der i dette „Givne“ findes
Forskelligheder og Sammenhænge, Kvaliteter og Kontinuiteter. James
fremhæver især de stadige Overgange (continuous transitions) og
Vanskeligheden ved at skelne mellem Forhold og Led (relation and
matter related); Bergson betoner især den indre Strøms usondrede og
dog uensartede Væsen (durée hétérogène et indistincte; durée dont les
moments se pénètrent). Men hermed er da angivet en indre Mulighed for
Eftertankens Opstaaen. Hele Vægten maa ikke lægges paa de ydre
Aarsager (det praktiske Behov, Sprogets Mekanik og Videnskabens
Teknik). Motivet til Eftertankens Vaagnen er givet med Modsætningen
mellem Kontinuitet og kvalitativ Mangfoldighed, mellem Forbindelsen og
de enkelte Elementer. Herved er der Mulighed for en Hæmning af den
uvilkaarlige Strøm, en Hæmning, hvorpaa kun den diskursive Eftertanke
kan raade Bod ved den Klarhed og Bestemthed, den formaar at bringe.
Som en Sammenfatten, der fuldbyrdes gennem Sammenligning, kan
Eftertanken tilvejebringe Klarhed over Forholdet mellem Kontinuitet og
Diskontinuitet, Enhed og Mangfoldighed. Og det er ikke noget helt
Fremmed, den bringer ind ved sit Arbejde; ti allerede i det
umiddelbart Givne fremtræde jo Ligheder og Forskelle, Sammenhæng og
Mangfoldighed --- og tillige, hvad der især er af Betydning, en
uvilkaarlig Stræben efter at bringe Harmoni tilstede mellem de
forskellige Tendenser. Hvis der ikke var en saadan oprindelig Stræben,
vilde det være uberettiget at tale om en Strøm, om en continuité
mouvante, en unité de direction\footnote{En af mine unge Tilhørere
skrev for nogle Aar siden til mig: „Jeg kan sige mig fri for ethvert
Formaal, ligesom for enhver Interesse eller Tro paa Noget udenfor
eller indeni mig. Jeg føler ganske vist noget Villende,
Fremadhastende, men det lever ligesom sit eget Liv, skubber Alt
tilside, der kommer ivejen, og begaar sikkert aldrig nogen Handling,
der ikke er til dets eget Bedste.“ --- Her har Eftertanken intet
Resultat bragt, kun Opløsning; men det uvilkaarlige Liv har holdt sig
i sin Enhed og i sin bestemte Retning, der for den unge Psykolog har
været saa meget des mere paafaldende, jo mindre hans Eftertanke har
kunnet naa til noget fast Stade.}! Der viser sig da her en Analogi til
det, vi for Eftertankens Vedkommende kaldte Syntese, en Tendens til at
samle og forene. Og Eftertanken vaagner, naar der er særlig Trang til
en fuldere Gennemførelse af denne Tendens, end muligt er indenfor det
uvilkaarlige Sjæleliv. Allerede indenfor dette, ikke først efterat
Eftertanken er vaagnet, kan der opstaa Kriser og Katastrofer og trues
med Kaos. Der kan kræves et Arbejde, allerede før Eftertanken kan
vaagne; men Vanskelighederne opstaa, bearbejdes og løses paa andre
Vilkaar i det uvilkaarlige Sjæleliv, end hvor Eftertanken har taget
Styret. Maaske foregaar noget af det Betydningsfuldeste i vort
Sjæleliv i Regionerne nedenunder Eftertankens Omraade, og det er ofte
kun Eftervirkningerne af disse dybtliggende Processer, Eftertanken faar
at gøre med. Det vil altid være umuligt at godtgøre, at vi paa noget
Punkt staa overfor et umiddelbart Givet i absolut Betydning, det vil
sige et saadant, hvor der aldeles intet psykisk Arbejde er gjort. Hvad
vi i hvert enkelt Tilfælde kalde umiddelbart givet, fortjener kun
dette Navn i Forhold til et vist bestemt psykisk Arbejde, medens det i
Sammenligning med andre Emner maaske vilde være at betragte som
Resultatet af et Arbejde.

4. Ved at gaa ud fra Eftertanken og søge at blive klar over dens Maade
at virke paa, kom vi til Begrebet Syntese, og vi fandt derefter, at
dette Begreb ogsaa kunde finde Anvendelse paa det uvilkaarlige
Sjæleliv, der danner det stadige Grundlag, af hvilket Eftertanken
udspringer for der-

% [text to be added: pp. 11--20]

\end{document}
